\clearpage
\fancyhead[L]{\small INSTRUMENTE.}
\hypertarget{native-427-part-lower}{}
\begin{center}
\Large INSTRUMENTE.
\end{center}



\begin{center}
Göttingische gelehrte Anzeigen. Stück 167. Seite 1665..1669. 1819. October 18.
\end{center}





So wie wir im 127. Stück dieser Blätter vom vorigen Jahr [S. 410 d. B.] eine Beschreibung des
\textsc{Repsold}schen Meridiankreises und einige der ersten damit auf der hiesigen Sternwarte angestellten
Beobachtungen mitgetheilt haben, geben wir jetzt eine kurze Nachricht von dem \textsc{Reichenbach}schen
\emph{Mittagsfernrohr}, als dem zweiten Hauptinstrumente, dessen Aufstellung seit kurzem vollendet ist. Da
die Einrichtung der Mittagsfernröhre, so wie sie bei dem gegenwärtigen Zustande der Beobachtungskunst
angewandt werden, in allen wesentlichen Beziehungen immer dieselbe ist, so dürfen wir uns hier auf
dasjenige beschränken, was dem hiesigen eigenthümlich ist. Das Fernrohr hat sechs Pariser
Fuss Brennweite und 52 Linien Oeffnung, und wird also in letzterer Rücksicht nur von dem neuen Mittagsfernrohr
der Greenwicher Sternwarte übertroffen, welches bei 10 Fuss Brennweite eine Oeffnung von
5 Engl. Zollen hat. Die Axe hält 37 Zoll, die Zapfen von Stahl ruhen in Pfannen von Glockenmetall.
Oculare sind 4, welche nach einer vorläufigen Bestimmung 80, 110, 150 und 210 mal vergrössern. Auch
die stärkste Vergrösserung gibt vollkommen scharfe Bilder. Hr. Hofr. \textsc{Gauss} bedient sich derselben gewöhnlich,
obwohl sie, da das Gesichtsfeld dabei nur 9 Minuten im Durchmesser hält, eine etwas sorgfältigere
Stellung der Höhe erfordert. Das Netz besteht aus Spinnenfäden; der verticalen sind sieben,
deren Zwischenräume von Sternen im Aequator in 10 Secunden durchlaufen werden. Hat man die Fertigkeit,
welche das Beobachten bei diesen kleinen Intervallen erfordert, einmal erworben, so findet man
es überaus angenehm, in kurzer Zeit eine so grosse Anzahl Antritte beobachten zu können, so wie es
bei dieser grössern Fädenzahl auch weniger auf sich hat, wenn, absichtlich oder zufällig, an einem
oder einigen Fäden kein Antritt beobachtet worden ist. Hr. Hofr. \textsc{Gauss} findet den sogenannten wahrscheinlichen
Fehler des beobachteten Antritts an einem Faden, aus mehreren hundert von ihm gemachten Beobachtungen von Sternen, die nicht zu weit vom Aequator
abstehen, $=0^{\mathrm s}095$. An der Libelle,
mit welcher die Axe horizontal gemacht wird, entspricht ein Ausschlag von einem Pariser Zoll einer
Neigung von 20 Secunden. Das Stellen des Fernrohrs für jede vorgeschriebene Declination geschieht
vermittelst eines kleinen am Fernrohr selbst nahe beim Ocularende befestigten, unmittelbar in Viertelsgrade
und durch den Vernier in Minuten getheilten Kreises, auf dem sich eine Alhidade mit einer
kleinen Libelle dreht; der Index an der Alhidade gibt sofort die Declination an. Diese Einrichtung
gewährt den grossen Vortheil, auch bei umgekehrter Lage des Instruments mit gleicher Leichtigkeit
auf jede Höhe richten zu können; durch einen zweiten ähnlichen auf der andern Seite anzubringenden
Kreis, welchen der Künstler noch nachliefern wird, wird noch der neue Vortheil erreicht werden, zwei
schnell nach einander durch die Mittagsfläche gehende Sterne ohne Zeitverlust beobachten zu können. ---
Die optische Vollkommenheit des Fernrohrs ist bewundernswürdig. Einige Proben, die wir hier aus den
Beobachtungen des Hrn. Hofr. \textsc{Gauss} anführen, mögen zeigen, was gute, jedoch nicht gerade ungewöhn%
\SourceLeaf{428}%
lich scharfe Augen, in unserm Clima bei günstigem aber doch keinesweges ausserordentlichem Zustande
der Luft, bei gehöriger Aufmerksamkeit, aber ohne peinliche Anstrengung, zu erkennen vermögen. Bei
\emph{voller Fadenbeleuchtung} ist ein Saturnstrabant, öfters zwei, gut zu erkennen und zu beobachten.
Eben so der Nebenstern des Nordsterns. Der kleine Stern, welcher jetzt nur 5 Minuten vom Pole absteht,
verträgt noch hinlängliche Beleuchtung, um gut beobachtet zu werden; schon in der Dämmerung, wo die
Fäden noch gar keiner Beleuchtung bedürfen, lässt er sich erkennen und beobachten. Bei Tage sind die
Sterne bis 3. Grösse, wenn sie weder zu tief, noch zu nahe bei der Sonne stehen, meistens mit
Leichtigkeit zu beobachten. Bei günstiger Luft und unter Anwendung einiger Vorsichtsregeln sind auch
Sterne vierter, ja selbst fünfter Grösse noch bei Tage recht gut zu erkennen. So hat z. B. Hr. Hofr.
\textsc{Gauss} die Sterne fünfter Grösse $\gamma'$ im kleinen Bär, $e$ im grossen Bär, $1$ im Drachen
nach Hevel, $g$ im grossen Bär (Alcor) bei hellem Tage, letztern selbst im Mittage beobachtet. Die
Schwierigkeit, dergleichen kleine Sterne bei Tage zu observiren, liegt eigentlich nur darin, dass die
Natur des Instruments kein ganz scharfes Einstellen verstattet: wo \emph{diese} Schwierigkeit wegfällt,
und das Auge also nur einen sehr kleinen Theil des Gesichtsfeldes zu fixiren braucht, erkennt man
Sterne 5ter Grösse, ja selbst noch kleinere, leicht. Dies ist der Fall bei Doppelsternen; z. B. der
Begleiter von $\zeta$ im grossen Bär, dem die Sternverzeichnisse nur die sechste Grösse beilegen, ist
am hellen Mittage sehr schön zu sehen; eben so ist der noch kleinere Begleiter von $\alpha$ im
Hercules, gegenwärtig, wo er Nachmittags um 4 Uhr culminirt, sehr gut zu erkennen. Uebrigens sind die
mannichfaltigen Vortheile, welche die Möglichkeit, auch kleinere Sterne bei Tage gut beobachten zu
können, gewährt, Kennern von selbst einleuchtend. Wir erwähnen hier nur des einen, sehr wichtigen,
dass man zur Berichtigung des Instruments durch Circumpolarsterne, die bekanntlich die einzige eines
solchen Instruments würdige, und so oft als möglich anzustellen ist, nicht auf den Nordstern
beschränkt bleibt, dessen eine Culmination oft in eine unbequeme Nachtstunde fällt, oder dessen
Beobachtung durch ungünstige Umstände vereitelt wird.

\par\indent Eine besondere Erwähnung verdient noch die Vorrichtung zum Ausheben und Umlegen des Instruments.
Aus den Lagern gehoben wird das Fernrohr vermittelst eines Seils, welches oben über eine Rolle geht,
und bis zu einer unten an der Mauer befestigten Winde geführt ist. Das Fernrohr wird dann frei in der
Luft schwebend ohne alle Erschütterung umgewandt, und wieder eingelegt. Die ganze Operation erfordert
nur einige Minuten, die über dem Fernrohre der Decke des Zimmers befindliche Rolle wird nach
vollendetem Geschäft, um beim Beobachten von Zenithalsternen nicht hinderlich zu sein, durch einen
Zug an die Seite geschoben. 
Dieser ganze schöne Apparat ist von dem hiesigen Hrn. Inspector \textsc{Rumpf} sehr zweckmässig und
solide ausgeführt.

\par\indent Wir fügen hier zur Probe noch einige von den ersten mit diesem Instrumente gemachten Beobachtungen
bei, und bemerken nur, dass die Declinationen diesmal nicht haben beobachtet werden können, weil der
\textsc{Repsold}sche Meridiankreis in diesem Monat abgenommen gewesen war.

\Needspace{27\baselineskip}
\begin{center}
\emph{Beobachtungen des Saturn.}

\small
\setlength{\tabcolsep}{5pt}
\begin{tabular}{l|rrr|rrr}
\multicolumn{1}{c|}{1819}  &  \multicolumn{3}{c|}{Mittlere Zeit}  &  \multicolumn{3}{c}{Scheinbare Gerade Aufsteig.}  \\
\hline
Sept.\ 4  &  $13^{\mathrm u}$  &  $5^{\mathrm m}$  &  $43^{\mathrm s}9$  &  $359^\circ$  &  $40'$  &  $2''7$  \\
5  &  $13$  &  $1$  &  $22.0$  &   &  $36$  &  $2.7$  \\
6  &  $12$  &  $57$  &  $10.7$  &   &  $32$  &  $9.6$  \\
7  &  $12$  &  $52$  &  $58.8$  &   &  $28$  &  $10.3$  \\
8  &  $12$  &  $48$  &  $46.1$  &   &  $24$  &  $4.0$  \\
9  &  $12$  &  $44$  &  $34.4$  &   &  $20$  &  $0.4$  \\
10  &  $12$  &  $40$  &  $22.2$  &   &  $15$  &  $55.6$  \\
11  &  $12$  &  $36$  &  $10.1$  &   &  $11$  &  $51.7$  \\
12  &  $12$  &  $31$  &  $57.5$  &   &  $7$  &  $41.2$  \\
\hline
 \multicolumn{1}{c|}{1819}  &  \multicolumn{3}{c|}{Mittlere Zeit}  &  \multicolumn{3}{c}{Scheinbare Gerade Aufsteig.}  \\
\hline
 Sept.\ 13  &  $12^{\mathrm u}$  &  $27^{\mathrm m}$  &  $45^{\mathrm s}0$  &  $359^\circ$  &  $3'$  &  $30''4$  \\
 14  &  $12$  &  $23$  &  $32.3$  &  $358$  &  $59$  &  $18.1$  \\
 15  &  $12$  &  $19$  &  $19.8$  &   &  $55$  &  $9.0$  \\
 16  &  $12$  &  $15$  &  $7.0$  &   &  $50$  &  $54.4$  \\
 19  &  $12$  &  $2$  &  $28.6$  &   &  $38$  &  $10.6$  \\
 21  &  $11$  &  $54$  &  $3.0$  &   &  $29$  &  $42.9$  \\
 24  &  $11$  &  $41$  &  $24.2$  &   &  $15$  &  $56.4$  \\
 28  &  $11$  &  $24$  &  $33.4$  &   &  $0$  &  $5.1$  \\
 Octob.\ 1  &  $11$  &  $11$  &  $55.9$  &  $357$  &  $47$  &  $36.3$  \\
\end{tabular}
\end{center}

\Needspace{22\baselineskip}
\SourceLeaf{429}
\emph{Beobachtungen der Vesta.}

\begin{center}
\begin{tabular}{l|rrr|rrr}
\multicolumn{1}{c|}{1819}  &  \multicolumn{3}{c|}{Mittlere Zeit}  &  \multicolumn{3}{c}{Scheinbare Gerade Aufsteig.}  \\
\hline
Sept.\ 7  &  $13^{\mathrm u}$  &  $32^{\mathrm m}$  &  $52^{\mathrm s}0$  &  $9^\circ$  &  $28'$  &  $7''2$  \\
9  &   &  $23$  &  $32.6$  &   &  $6$  &  $10.5$  \\
10  &   &  $18$  &  $51.1$  &  $8$  &  $54$  &  $44.7$  \\
11  &   &  $14$  &  $8.6$  &   &  $43$  &  $1.8$  \\
12  &   &  $9$  &  $25.1$  &   &  $31$  &  $6.9$  \\
15  &  $12$  &  $55$  &  $18.7$  &  $7$  &  $53$  &  $49.8$  \\
\hline
 \multicolumn{1}{c|}{1819}  &  \multicolumn{3}{c|}{Mittlere Zeit}  &  \multicolumn{3}{c}{Scheinbare Gerade Aufsteig.}  \\
\hline
 Sept.\ 16  &  $12^{\mathrm u}$  &  $50^{\mathrm m}$  &  $21^{\mathrm s}4$  &  $7^\circ$  &  $40'$  &  $58''5$  \\
 19  &   &  $35$  &  $55.5$  &   &  $1$  &  $18.0$  \\
 21  &   &  $26$  &  $15.9$  &  $6$  &  $34$  &  $15.9$  \\
 24  &   &  $11$  &  $43.4$  &  $5$  &  $52$  &  $58.0$  \\
 28  &  $11$  &  $52$  &  $18.1$  &  $4$  &  $57$  &  $24.3$  \\
 Octob.\ 1  &   &  $37$  &  $45.3$  &   &  $16$  &  $1.2$  \\
\end{tabular}
\end{center}

\Needspace{27\baselineskip}
\par\indent Da diese Beobachtungen noch nicht mit Tafeln oder Elementen verglichen sind, so setzen wir auch noch die Beobachtungen eines Fixsterns her, nebst der Reduction auf den mittlern Ort für den Anfang des Jahrs 1820.


\begin{center}
\emph{Gerade Aufsteigung von 85 Pegasi.}
\end{center}

\begin{center}
\begin{tabular}{l|rrr|rrr}
\multicolumn{1}{c|}{1819}  &  \multicolumn{3}{c|}{Beobachtete}  &  \multicolumn{3}{c}{Mittlere für 1820}  \\
\hline
Sept.\ 4  &  $358^\circ$  &  $11'$  &  $40''8$  &  $358^\circ$  &  $11'$  &  $43''0$  \\
5  &   &   &  $41.4$  &   &   &  $43.4$  \\
6  &   &   &  $40.2$  &   &   &  $42.0$  \\
7  &   &   &  $43.9$  &   &   &  $45.5$  \\
8  &   &   &  $42.6$  &   &   &  $44.0$  \\
9  &   &   &  $42.9$  &   &   &  $44.1$  \\
10  &   &   &  $43.0$  &   &   &  $44.1$  \\
\hline
 \multicolumn{1}{c|}{1819}  &  \multicolumn{3}{c|}{Beobachtete}  &  \multicolumn{3}{c}{Mittlere für 1820}  \\
\hline
 Sept.\ 11  &  $358^\circ$  &  $11'$  &  $42''6$  &  $358^\circ$  &  $11'$  &  $43''6$  \\
 12  &   &   &  $45.4$  &   &   &  $46.2$  \\
 13  &   &   &  $41.8$  &   &   &  $42.4$  \\
 14  &   &   &  $44.7$  &   &   &  $45.1$  \\
 15  &   &   &  $44.8$  &   &   &  $45.1$  \\
 16  &   &   &  $43.0$  &   &   &  $43.2$  \\
\multicolumn{4}{r|}{Im Mittel . .} &  $358$  &  $11$  &  $43.97$  \\
\end{tabular}
\end{center}

\par\indent Dieser Stern zeichnet sich durch eine starke eigene Bewegung aus; die Vergleichung seiner aus diesen Beobachtungen hervorgehenden geraden Aufsteigung mit der \textsc{Bradley}schen für 1755 gibt die eigne Bewegung in 64.7 Jahren $+54''65$, also in Einem Jahre $+0''845$.
